% =====================================================================
% §4 — Il sistema nella tradizione
%
% Struttura: genealogia in tre atti (differito per necessità ->
%            svolta real-time documentata nel volume X 1993 ->
%            ritorno volontario al differito) -> paragrafo onesto
%            «che cosa non è nuovo» (con i precursori nominati da noi)
%            -> le tre proposte, dimensionate, ciascuna col proprio
%            precursore più vicino -> chiusura che consegna a §5
%
% TODO prima della submission:
%   [x] chiavi BibTeX: FATTO (2026-06-12) — allineate alle chiavi reali di
%       refs.bib: Sparano2018cim->Sparano2018, DeTintis1995cim->DeTintis1995,
%       RolfeKeller2000cim->RolfeKeller2000, ArcellaSilvestri2012cim->Arcella2012,
%       Solomos2003ent->Solomos2003 (entretien citato col libro, pagine nel testo)
%   [x] verifica letteratura per il gate: FATTA (2026-06-11) su AC Toolbox,
%       Common Music, CMask — nessun gate di probabilità per-evento come
%       parametro di prima classe; la rivendicazione è stata circoscritta
%       di conseguenza. CHIUSO (2026-06-12): ingest CMask nel wiki
%       (sources/papers/bartetzki1997.md + concepts/deviazione-ampiezza-
%       probabilita.md) + entry refs.bib Bartetzki1997 completata (articolo,
%       STEAM/HfM Hanns Eisler Berlin). Verificato sul manuale: il quantizer CMask
%       ha una strength continua (attrazione per-valore, envelope-abile), non
%       un gate Bernoulli per-evento — la distinzione regge la rivendicazione.
%   [ ] budget citazionale complessivo del paper: target 9–21 ref;
%       valvole di taglio segnalate con %VALVOLA
% =====================================================================

% ================================================================
\section{Il sistema nella tradizione}\label{sec:tradizione}

Storicamente, il tempo differito della sintesi granulare non è una
scelta: è la condizione di partenza. La prima implementazione
documentata~\cite{Roads1978} e il primo articolo CIM dedicato alla
tecnica~\cite{Roads1985cim} formulano il problema del controllo ---
migliaia di eventi al secondo, intrattabili uno a uno --- quando il
calcolo in tempo reale semplicemente non è disponibile; ed è già
lì, in tempo differito, che compare il pattern del front-end
dichiarativo che pilota un engine di sintesi. La svolta è
documentabile con precisione negli atti: nel volume del~1993 coesistono
la granulazione offline su mainframe di Di Scipio e
Tisato~\cite{DiScipioTisato1993cim} e il \emph{granular sampling} in
tempo reale su ISPW di Lippe~\cite{Lippe1993cim} --- lo snodo di
paradigma osservato nel suo farsi. Da allora il tempo reale è divenuto
la norma, in una linea che in ambito CIM arriva a
GrainLab~\cite{Sparano2018} e, fuori, a
EmissionControl2~\cite{Roads2021}. Il sistema qui descritto torna al
tempo differito per scelta: questa sezione ne traccia la genealogia, la
prossima ne discute le ragioni.

Quasi nulla dei meccanismi di \S\ref{sec:architettura} è nuovo, ed è
bene dirlo con precisione. La tendency mask è di
Truax~\cite{Truax1988}, e nel giro di pochi anni è nomenclatura
canonica anche nei Colloqui: la adottano nel~1993 tanto il sistema
offline di Di Scipio e Tisato (campionamento gaussiano su maschere)
quanto quello real-time di Lippe (\emph{«constantly moving windows with
varying sizes»}), e nel~1995 De Tintis la cita come stato
dell'arte~\cite{DeTintis1995}: il modello che il sistema eredita è
una postura consolidata della tradizione, non un'invenzione. Il modello
sincrono/asincrono della griglia è dello stesso Truax. E il pattern del
front-end dichiarativo che genera lo score granulare, anticipato in
ambito CIM da Roads, è realizzato compiutamente da CMask di
Bartetzki~\cite{Bartetzki1997}: maschere di tendenza per ogni campo
dello score, bordi mossi nel tempo da funzioni a segmenti, Csound in
uscita. La
micro-deviazione per grano ha radici altrettanto profonde: già Roads
nel~1985 ne quantifica l'effetto spettrale, e la decorrelazione come
proprietà della massa granulare è teorizzata in ambito CIM da Rolfe e
Keller~\cite{RolfeKeller2000} e, sul piano compositivo, da
Vaggione~\cite{Vaggione2002}. Perfino la famiglia di controllo è una
scelta dentro un ventaglio noto: Di Scipio governa gli stessi parametri
con mappe caotiche~\cite{DiScipio1991cim}, e il sistema affianca quella
famiglia senza pretendere di sostituirla.

Su questo fondo, le proposte del sistema sono tre, e vanno dimensionate.
La prima è lo YAML come notazione: una specifica dichiarativa completa,
validata durante la scrittura, che è insieme il documento di lavoro e
l'oggetto che si spedisce; dentro questo modello di controllo, la
fattorizzazione della deviazione per grano in ampiezza e probabilità
indipendenti, entrambe componibili come envelope
(\S\ref{sec:deviazione}), circoscrive la proposta: nei generatori di
score di questa famiglia il valore è sempre estratto dalla maschera, e
negli ambienti Lisp un gate è costruibile come idioma ma non è
parametro di prima classe del modello; come asse dichiarativo, per
quanto ci risulta, non ha precedente diretto. La seconda è la partitura
con l'asse verticale sulla posizione di lettura: il precursore concreto
è l'overlay multi-parametro su terminale di Truax (1988, Fig.~4), e la
descrizione verbale del meccanismo --- il movimento della testina nel
buffer rispetto al tempo di uscita --- è in
Truax~\cite{Truax1994}; la partitura lo rende disegno, là dove la
notazione granulare di Roads vive sul piano tempo--frequenza. La terza è
il workflow per stem con cache incrementale ed export DAW
(\S\ref{sec:render}): il suo parente in tempo reale è l'aspetto
«ricorsivo» di Lippe --- rimissaggio e riuso dell'uscita granulare ---
e quello compositivo il montaggio multitraccia praticato da Vaggione,
qui automatizzato e ancorato a un fingerprint per stream.

Nessuna delle tre proposte tocca la sintesi; tutte toccano il punto in
cui la tradizione colloca da sempre la difficoltà, il controllo. Ed è
per questo che il tempo differito non è, qui, un dettaglio
implementativo: è il presupposto che le rende possibili tutte e tre.
Resta da argomentare perché valga la pena sceglierlo oggi, mentre il
tempo reale è disponibile.


